\subsection{Agent-based modelling and simulation}
\label{sec:abms-netlogo}
\slidesrc{ABMS: 2--4, 18--21}

\begin{nbdef}{Agent-based modelling and simulation (slide 18)}
\begin{itemize}[nosep]
  \item \textbf{Agent-based modelling} is a computational methodology that allows scientists
        to create, analyse and experiment with \textbf{artificial worlds} populated by agents
        that interact in non-trivial ways and that constitute their own environment.
  \item An \textbf{agent-based simulation} (ABS) is one in which a collection of agents
        compete or collaborate, or both, to attain individual goals in pursuit of social or
        organisational goals.
  \item A simulation based on software agents can support good science provided that the
        design of the agents and of the simulation environment are themselves based on good
        science.
\end{itemize}
\end{nbdef}

\begin{nbfact}{Why ABMS? (slide 19)}
\begin{itemize}[nosep]
  \item Aspects of the real world we want to examine are \textbf{not physically
        accessible}; ABMS takes place in an artificial world.
  \item Experimenting with the real system is \textbf{prohibited} due to undesirable
        disturbances.
  \item The \textbf{time scale} of the system's behaviour is too small or too large for
        observation.
  \item The original system \textbf{does not exist any more, or does not exist yet}.
  \item Modelling is a tool for understanding: it \textbf{formalises a hypothesis} that
        otherwise would have remained very vague.
\end{itemize}
\end{nbfact}

\begin{center}
\small
\begin{tabularx}{\linewidth}{@{}l X@{}}
\toprule
\textbf{Reason} & \textbf{Example (added)} \\
\midrule
not accessible & the spread of a virus inside a city: we cannot observe every contact \\
must not be disturbed & testing an evacuation plan by setting a real stadium on fire \\
wrong time scale & climate and land-use change over a century; molecules in nanoseconds \\
does not exist (yet / any more) & traffic in a city district not yet built; the fall of an ancient civilisation \\
formalising a hypothesis & ``segregation can arise even if nobody wants it'' (Schelling, below) \\
\bottomrule
\end{tabularx}
\captionof{table}{The reasons of slide 19 with examples.}
\end{center}

\paragraph{Applications (slide 20).} Biology, energy, transportation, entertainment, and
many more (a Google Scholar search on applications of ABMS since 2020 returns thousands of
papers). The slides cite \textbf{MASSIVE}, a commercial agent-based crowd-simulation software
used to animate large battle and crowd scenes in films, where each digital character is an
agent with its own perception and behaviour.

\subsubsection{NetLogo}
\slidesrc{ABMS: 4, 21}

\textbf{NetLogo} (Northwestern University) is a free programmable modelling environment for
ABMS. Its world is a grid of \emph{patches}; mobile agents are called \emph{turtles}; an
\emph{observer} gives orders. It comes with a large \emph{models library} (slide 4: biology,
chemistry and physics, computer science, earth science, games, mathematics, social science;
slide 21 shows thumbnails of many models).

\begin{nbexample}{Schelling's segregation model (the model selected on slide 4)}
The model on slide 4: two types of turtles, red and green, live in a ``mythical pond''. They
get along with one another, but each turtle wants to make sure that it lives near
\emph{some} of ``its own'': each red turtle wants at least some red neighbours, each green
turtle at least some green ones. The simulation shows how these individual preferences
ripple through the pond, leading to large-scale patterns. (It is inspired by Thomas
Schelling's writings about social systems such as housing patterns in cities.)

The rule of each agent, at each tick (simplified sketch in NetLogo style):
\begin{minted}[linenos=false]{text}
to go
  ask turtles [
    let total   count (turtles-on neighbors)
    let similar count (turtles-on neighbors) with [color = [color] of myself]
    if total > 0 and similar / total < 0.3 [ move-to one-of patches with [not any? turtles-here] ]
  ]
  tick
end
\end{minted}
Each turtle is \emph{happy} if at least 30\% of its neighbours (the 8 surrounding cells) have
its colour; otherwise it moves to a random empty cell. Example: a red turtle with neighbours
\{green, green, red, green, empty, \ldots\}: \code{total} $= 4$, \code{similar} $= 1$,
$1/4 = 25\% < 30\%$: it moves. With 2 red out of 4, $50\% \geq 30\%$: it stays.

\textbf{What emerges.} Starting from a random mix, after a few dozen ticks the pond is
divided into large red and green regions, with typically far more than 30\% similar
neighbours on average. Nobody wanted segregation (everyone accepts being a minority of 70\%
in its neighbourhood), yet the city is segregated: a global pattern that no individual
intends, visible only by running the simulation. This is the ``third way of doing science''
of Section~\ref{sec:uses-of-mas}.
\end{nbexample}

\begin{nbintuition}[Simple agents, complex worlds]
The agents of an ABMS are usually \emph{reactive} (condition--action rules,
Section~\ref{sec:reactive-architectures}): that is why the professor listed ABMS among the
pros of reactive architectures (Whiteboard~\ref{wb:reactive-pros-cons}). The complexity is
not inside each agent, but in their \emph{interaction} (Section~\ref{sec:what-is-a-mas}).
In contrast, the BDI agents of Section~\ref{sec:bdi-agentspeak} are \emph{cognitive}: richer
individuals, typically fewer of them, used when the intelligence of the individual matters.
\end{nbintuition}
